\subsection{The first obstructions and the exact additive defect}

\begin{theorem}[E.15 --- Proved: the polarization law for the actual class-two obstruction]
\label{E-polarization}
Let $\ell$ be odd, let $N$ be a free pro-$\ell$ group of
finite rank with a specified continuous action of a
profinite group $G$, and put
$A=N^{\mathrm{ab}}$, $B=\bigwedge^2 A$.
A choice of free lifts of a basis of $A$ gives the integral
class-two coordinates
\[
 N/\gamma_3N=A\times B,\qquad
 (v,m)(w,n)=(v+w,m+n+\tfrac12v\wedge w).
 \tag{E15.1}
\]
Write the genuine arithmetic action as
\[
 g(v,m)=(A_gv,\ B_gm+D_gv),\qquad
 B_g=\bigwedge^2 A_g.
\]
For continuous one-cocycles $a,b:G\to A$, the obstruction
to lifting $a$ through the actual class-two extension is
represented by
\[
 \omega_a(g,h)=D_g a(h)+\tfrac12a(g)\wedge A_g a(h).
 \tag{E15.2}
\]
Its exact polarization is
\[
 \operatorname{ob}(a+b)-\operatorname{ob}(a)
       -\operatorname{ob}(b)=[a\cup_{\wedge}b]
       \quad\text{in }H^2_{\rm cts}(G,B).
 \tag{E15.3}
\]
Thus the sum of two liftable classes is liftable precisely
on the zero locus of their displayed cup pairing.

For the actual quartic decomposition of E.17,
$A=W\oplus j(V^3)$ at
$\ell=3,5$, where $V=T_\ell E_+$, the mixed component
for $a=b+j_{\mathbf r}h$ has at every good prime
$p\ne2,\ell$ the following actual local representative.
Let $v_p=b(t_p)$ and
$u_p=(F_{E,p}-1)^{-1}h(F_p)$, with the choices of E.14.
Its unique rational tame correction is
\[
 c_p=v_p\wedge j_{\mathbf r}u_p .
 \tag{E15.4}
\]
Writing $t=\min_i v_\ell(r_i)$ and
$\operatorname{ord}(u_p+V)=\ell^{m_p}$, this correction
is integral exactly when
\[
 v_\ell(v_p)+t\ge m_p\qquad(m_p>0).
 \tag{E15.5}
\]
After replacing $u_p$ by its fractional representative,
the sequence of corrections is restricted exactly when
$v_\ell(v_p)+t-m_p\to+\infty$ on the primes with $m_p>0$.
These are the polarization and mixed support maps in
LM7 E.115 and E.117, with the arithmetic linear term
retained. The density of the local inequality's failure
at a fixed finite threshold $d=v_\ell(v_p)+t$ is the
tail $\delta_\ell(m_p\ge d+1)$ of E.14; this statement
concerns the local mixed test with that fixed threshold.
\end{theorem}
\begin{proof}
The Hall commutator isomorphism
$\gamma_2N/\gamma_3N\simeq\bigwedge^2A$ gives (E15.1).
Division by two is allowed in $\mathbb Z_\ell$.
An automorphism of this group is determined by its
linear action on $A$ and its remaining central correction.
Preservation of the multiplication shows that the
correction $D_g$ is additive. Applying the action
composition to $(v,0)$ gives
$D_{gh}=B_gD_h+D_gA_h$.
The section lift $g\mapsto(a(g),0,g)$ has multiplication
defect (E15.2); associativity makes it a two-cocycle.
Changing a continuous lift adds a continuous coboundary.
This proves the obstruction assertion for the genuine
action, without setting $D_g$ to zero.

At the cochain level the left side of (E15.3) is
\[
 \tfrac12\{a(g)\wedge A_gb(h)+b(g)\wedge A_ga(h)\}.
\]
Set $q(g)=a(g)\wedge b(g)$ and use
$dq(g,h)=B_gq(h)-q(gh)+q(g)$. The two one-cocycle
identities give
\[
 dq(g,h)=-a(g)\wedge A_gb(h)+b(g)\wedge A_ga(h).
\]
The preceding polarization is therefore
$a\cup_\wedge b+\tfrac12dq$, proving (E15.3).
Its vanishing equivalence follows by lifting the
central extension. In particular this formula
specifies the entire additive defect, including
the possible mixed terms of two general inputs.

For the local mixed formula, consider the central
extension represented by the ordered cup
$\omega(g,h)=b(g)\wedge g e(h)$, where
$e=j_{\mathbf r}h$ is unramified at $p$.
In its multiplication
$(z,g)(w,h)=(z+gw+\omega(g,h),gh)$, choose central-zero
lifts of tame inertia. The relation
$F_pt_pF_p^{-1}=t_p^p$ has central defect
\[
 -p\,v_p\wedge e(F_p).
\]
Changing the inertia lift by $z$ changes the defect by
$(F_{B,p}-p)z$. The toric action is $F_{W,p}=pP_p$,
where $P_p$ is the boundary permutation. The
one-cocycle tame relation gives $P_pv_p=v_p$.
Consequently
\[
 (p^{-1}F_{B,p}-1)
       (v_p\wedge j_{\mathbf r}u)
   =v_p\wedge j_{\mathbf r}(F_{E,p}-1)u.
\]
Every eigenvalue of this mixed module has complex
absolute value $p^{3/2}$, so the displayed operator
is invertible over $\mathbb Q_\ell$. This proves
(E15.4) as its unique rational solution.

The natural map $W\otimes j(V^3)\to W\wedge j(V^3)$
is an integral lattice isomorphism, because these
are complementary direct summands of $A$.
A nonzero pure tensor has valuation equal to the
sum of the two vector valuations: choose primitive
coordinates in each factor and inspect their
product coordinate. A fractional representative of
$u_p+V$ is zero for $m_p=0$ and has valuation $-m_p$
otherwise. This proves (E15.5) and the limit statement.
Replacing a representative by an integral vector
changes $c_p$ by an integral restricted family,
because the inertia coordinates of a continuous
global cocycle tend to zero. E.14 supplies the
claimed density at every fixed threshold.
The passage is \S51.16, folios 733v--736v,
scans 172--178, and the explicit quartic task E.23(iv).
\end{proof}

\begin{theorem}[E.16 --- Proved: full pro-three and pro-five sections beyond every nilpotent stage]
\label{E-regular-prime-family}
Let $\ell\in\{3,5\}$, $r\ge1$, and $n=\ell^r$.
Use the rational unit tangent at
$[0:-1:1]\in X_n$ and the resulting genuine action
on
\[
 N_n=\pi_1(U_{n,\overline{\mathbb Q}})^{(\ell)},\qquad
 V_n=N_n^{\rm ab},\qquad S=\{\ell,\infty\}.
\]
Every class of $H^1(G_{\mathbb Q,S},V_n)$ lifts to a
section of the full geometrically pro-$\ell$
arithmetic extension of $U_n/\mathbb Q$.
There are infinitely many such classes whose lifts
are cuspidal at no boundary point. Each retains
this property in every positive-class intrinsic
nilpotent quotient and every positive-length
solvable quotient, through the compactified
Jacobian quotient specified below.
Here the cuspidal classes in a quotient are the
images of the actual cusp packets under the
displayed arithmetic quotient map.

Thus no intrinsic pro-$3$ nilpotent level for
$n=3^r$, and no intrinsic pro-$5$ nilpotent level
for $n=5^r$, selects only cuspidal sections. This
extends the checked mechanism of LM7 E.34 from
the cubic to these two entire towers. It gives
a specific range of the full section problem;
the passage to all coefficient primes and
mixed finite quotients is retained in E.21.
\end{theorem}
\begin{proof}
Put $K=\mathbb Q(\zeta_\ell)$ and let $S_K$ contain
the unique prime above $\ell$. Both fields have
class number one. For $\ell=3$ the Eisenstein
integer norm is Euclidean. For $\ell=5$, the
degree-four field has discriminant $125$, and
Minkowski's ideal-class bound is
\[
 \frac{4!}{4^4}\left(\frac4\pi\right)^2\sqrt{125}<2.
\]
Every ideal class therefore contains an integral
ideal of norm one. The factorization
$\ell=\text{unit}\cdot(1-\zeta_\ell)^{\ell-1}$
gives the single prime in $S_K$.

Write $r_2=(\ell-1)/2$. Unique factorization,
the $S_K$-unit theorem and Kummer theory give
\[
 \dim_{\mathbb F_\ell}
 H^1(G_{K,S_K},\mathbb F_\ell)=r_2+1 .
\]
Indeed the ordinary units contribute rank $r_2-1$,
the roots of unity contribute one dimension, and
the $S_K$-valuation contributes one. The global
Euler formula, in its totally imaginary case, is
\[
 \frac{\#H^0\#H^2}{\#H^1}=\ell^{-r_2}.
\]
Since $\#H^0=\ell$, this proves
$H^2(G_{K,S_K},\mathbb F_\ell)=0$.
The finiteness and Euler inputs are the exact
I.4.10, I.4.15 and I.5.1 statements checked in
E.13~\cite{E_density_MilneADT}.

The maximal pro-$\ell$ quotient $P$ of this Galois
group is free of rank $r_2+1$. To check the
cohomological criterion here, let $J$ be the kernel
of the maximal quotient. A nonzero invariant
character $J\to\mathbb F_\ell$ would produce a
pro-$\ell$ extension of $P$ and contradict its
maximality. Inflation therefore injects
$H^2(P,\mathbb F_\ell)$ into the zero group above.
For a finitely generated pro-$\ell$ group, this
vanishing solves every finite embedding problem
with $\ell$-group kernel: filter an abelian kernel
by trivial simple modules, and then induct on
central characteristic subgroups of a general
finite kernel. A minimal free presentation splits
by compactness of its finite solution sets; its
idempotent splitting is the identity on the
Frattini quotient and hence on the free group.
This proves freeness with the stated rank.

Let $F$ be the free pro-$\ell$ tripod group, with
its standard rational tangent. The actual power
cover identifies
\[
 N_n=\ker(F\longrightarrow(\mathbb Z/n)^2).
 \tag{E16.1}
\]
At the chosen cusp, the equations are
$t=s^n$, $y/z=-(1-s^n)^{1/n}$.
Their leading constants have coherent rational
$\ell$-power roots. Coefficientwise action on
these Puiseux fibre functors therefore gives
the genuine restriction of the tripod action
in (E16.1). Smooth marked specialization away
from $\ell$, with this integral unit tangent,
makes that action unramified there.

Over $K$ the action on $F/\Phi(F)$ is trivial,
since $\mu_\ell\subset K$. The kernel of
$\operatorname{Aut}(F)\to\operatorname{GL}(F/\Phi(F))$
is pro-$\ell$: on the characteristic lower
exponent-$\ell$ quotients its successive kernels
embed in homomorphism groups of finite
$\mathbb F_\ell$-spaces, and the finite basis
theorem gives the initial step. Hence the
action on $F$ and on $N_n$ factors through $P$.
Over $\mathbb Q$ it factors through the actual
relative quotient
\[
 1\longrightarrow P\longrightarrow\Gamma
 \longrightarrow\operatorname{Gal}(K/\mathbb Q)
 \longrightarrow1 .
\]
The quotient has order $\ell-1$. Every finite
$\ell$-primary module through $\Gamma$ has zero
second cohomology over $G_{\mathbb Q,S}$:
over $G_{K,S_K}$ filter it by trivial
$\mathbb F_\ell$ factors and use the preceding
vanishing; then use the exact averaging
operator for the quotient of order $\ell-1$.

Every continuous $V_n$-cocycle factors through
this same $\Gamma$. On $G_{K,S_K}$ its graph
takes values in the pro-$\ell$ group
$V_n\rtimes\rho(G_{K,S_K})$ and therefore kills
$J$. Finite embedding problems now lift the
graph through every finite characteristic
quotient of $N_n$. There are finitely many
solutions at each finite level, since
$\Gamma$ is finitely generated; any finite
set of level constraints is solved at a
common refinement. Compactness gives
compatible solutions and hence the asserted
continuous full lift.

To produce noncuspidal lifts, use the integral
exact sequence from the generalized Jacobian:
\[
 0\longrightarrow W_n\longrightarrow V_n
 \longrightarrow M_n=T_\ell\operatorname{Jac}(X_n)
 \longrightarrow0 .
 \tag{E16.2}
\]
The toric term $W_n$ has the same pro-$\ell$
action factorization. The finite-coefficient
vanishing above and finite $H^1$ groups give
$H^2_{\rm cts}(G_{\mathbb Q,S},W_n)=0$.
Thus (E16.2) surjects on $H^1$.
The group $H^1(G_{\mathbb Q,S},M_n)$ is finitely
generated over $\mathbb Z_\ell$. Its rank is
\[
 g(X_n)=\frac{(n-1)(n-2)}2>0 .
\]
Here $H^2_{\rm cts}$ vanishes by the same argument,
$H^0(M_n)=0$ by Mordell--Weil, and complex
conjugation has $g(X_n)$ eigenvalues of each
sign. The integral Euler formula gives the
displayed rank; possible finite torsion in
$H^1$ does not affect it. Choose any infinite-order
class, lift it through (E16.2), and then apply
the full lifting construction.

For any two Fermat cusps $d,O$, the ratio of
their tangent-line equations has divisor
$n(d-O)$. Each tangent line meets the Fermat
curve only at its tangent cusp, with
multiplicity $n$. Thus every cusp class in
the proper Jacobian is $n$-torsion.
The compactified image of a cuspidal section
is the Kummer class of that cusp; changing its
tangent affects only $W_n$. Our chosen
infinite-order compact class cannot be such
a class. Distinct geometric cusp inertia
classes stay distinct in the free pro-$\ell$
group, so a section normalizing one requires
the underlying cusp to be rational. There
are exactly three rational cusps for odd $n$.
The same compactified map factors through
every quotient named in the theorem, and
retains the obstruction to cuspidality.
Varying the infinite-order class proves
infinitude. We claim no priority for the
general pro-prime phenomenon; compare the
published constructions of Hoshi~\cite{Hoshi2010}.
The manuscript comparison is \S51.16,
folios 733v--736v, scans 172--178.
\end{proof}

\begin{proposition}[E.17 --- Proved: the actual quartic coefficient maps and a class-two origin]
\label{E-quartic-origin}
Put $C:x^4+y^4+z^4=0$, $U=C\setminus\{xyz=0\}$,
$F=\mathbb Q(\zeta_8)$, and $E_+:Y^2Z=X^3+XZ^2$.
For every odd prime $\ell$, the actual arithmetic
extension with geometric kernel
$N/\gamma_3N$, $N=\pi_1(U_{\overline{\mathbb Q}})^{(\ell)}$,
has a continuous section $S_\ell$. Its image in the
rank-two coordinate-unit extension is the zero section.
There is a specified $G_{\mathbb Q}$-equivariant integral
direct sum
\[
 N^{\rm ab}=W\oplus j((T_\ell E_+)^3),
 \qquad \rank W=11,
 \tag{E17.1}
\]
with the forward and inverse maps below. Relative to
$S_\ell$ the formulas of E.15 therefore apply to the
genuine quartic extension. These assertions verify the
precise steps of LM7 E.103 and E.105 used here; no
rank-one assertion about elliptic global cohomology is
needed for this construction or for E.13--E.15.
\end{proposition}
\begin{proof}
Let $J=\Jac(C)$ and $G=J_{C,D}$, where $D=C\setminus U$.
Write $p:G\to J$. A tangent-line ratio at two cusps has
divisor four times their difference. The extension
$1\to T_D\to G\to J\to1$ is classified by these cusp
differences in the dual Jacobian. Pullback by $[4]$
therefore splits. Its splitting is unique because
$\Hom(J,T_D)=0$; hence it descends to $\mathbb Q$ and
gives $j_0:J\to G$ with $pj_0=[4]$.

The three morphisms
\[
 \begin{aligned}
 f_0[x:y:z]&=[-x^2y:z^2x:y^3],\\
 f_1[x:y:z]&=[-y^2z:x^2y:z^3],\\
 f_2[x:y:z]&=[-z^2x:y^2z:x^3]
 \end{aligned}
 \tag{E17.2}
\]
have no simultaneous zero in their coordinate triples
on $C$, and direct substitution gives the equation of
$E_+$. For $f_0$, coordinates $u=x/y$, $v=z/y$ give
$X=-u^2$, $Y=uv^2$, so a generic fibre has four points.
The cyclic formulas have the same degree. Their three
different sign characters for the coordinate group
$C_2^2$ give $f_{i*}f_j^*=0$ for $i\ne j$: such a
composite equals its negative, and a homomorphism of
abelian varieties has no nonzero torsion. The diagonal
composites are $[4]$ by push--pull. Consequently
\[
 \Phi=(f_{i*}):J\to E_+^3,\quad
 \Psi=\sum_i f_i^*:E_+^3\to J,\quad
 \Phi\Psi=\Psi\Phi=[4].
\]
The second equality follows also because the image of
$\Psi\Phi-[4]$ is connected and contained in the finite
kernel of the isogeny $\Phi$. The map denoted $j$ in
(E17.1) is $j_{0*}\Psi_*/16$. The actual inverse maps are
\[
 a\longmapsto
 \bigl(a-j_{0*}p_*a/4,\ \Phi_*p_*a\bigr),
 \qquad
 (w,e)\longmapsto w+j_{0*}\Psi_*e/16.
 \tag{E17.3}
\]
All divisions are integral for odd $\ell$.

We verify the origin at the level of actual paths.
Set $\alpha_i=\zeta_8^{1+2i}$ and use the following
points, local parameters, and tangent lines:
\[
 \begin{array}{c|c|c|c}
 d&\text{point}&s_d&L_d\\\hline
 A_i&[0:\alpha_i:1]&x/z&y-\alpha_i z\\
 B_i&[\alpha_i:0:1]&y/z&x-\alpha_i z\\
 C_i&[\alpha_i:1:0]&z/y&x-\alpha_i y
 \end{array}
 \quad O=B_0 .
 \tag{E17.4}
\]
All tangent vectors have parameter value one. The
nonzero varying coordinate is
$\alpha_i(1+s_d^4)^{1/4}$; thus the leading coefficient
of $L_d$ is $\alpha_i/4$ in order four. Frame
$z_d=\mathcal O(d-O)\in G(F)$ by $s_d^{-1}$ at $d$,
$s_O$ at $O$, and its canonical section elsewhere.
For $f_d=L_d/L_O$, the torus coordinates of $4z_d$ are
\[
 t_{q,d}=
 \frac{\operatorname{lead}_q f_d}{
       \operatorname{lead}_O f_d}
       =\zeta_8^{K_{qd}}(\sqrt2)^{N_{qd}}.
 \tag{E17.5}
\]
This is a directly evaluable formula: the denominator
is $4L_d(O)/\alpha_0$; the numerator is
$\alpha_d/(4L_O(d))$ for $q=d$ and $L_d(q)/L_O(q)$
otherwise. The independently executed program
\texttt{quartic\_frames.py} evaluates all $121$ entries
in $\mathbb Q[z]/(z^4+1)$ using rational coefficients,
and records both $N$ and the phases $K$. With the order
$A_0,\ldots,A_3,B_1,B_2,B_3,C_0,\ldots,C_3$,
the three columns used below have average
\[
 \frac1{32}\sum_{d=B_1,B_2,B_3}N_{\bullet d}
 =v=\left(-\tfrac12,-\tfrac12,-\tfrac12,-\tfrac12,
 -\tfrac58,-\tfrac34,-\tfrac58,
 -\tfrac12,-\tfrac12,-\tfrac12,-\tfrac12\right).
 \tag{E17.6}
\]
The entry formulas prove the computation without a
numerical field approximation.

Put $H=G_F$ and choose coherent odd roots of $2$.
Choose roots of $\sqrt2$ whose squares are these roots;
then $\kappa_{\sqrt2}=\kappa_2/2$. All coherent odd
roots of $\mu_8$ lie in $F$, so their $H$-cocycles
vanish. For an odd power $m$ of $\ell$, choose an
$m$-division $v_m$ of $4z_d=t_d$. Integers $a,b$ with
$4a+mb=1$ give $u_m=a v_m+bz_d$, satisfying
$mu_m=z_d$, $4u_m=v_m$. These two equations determine
$u_m$ uniquely and compatibly, because $[4]$ is an
automorphism on $G[m]$. Its proper image is the fixed
four-torsion point $b[d-O]$. Thus the actual abelian
path cocycle, in (E17.3), is
\[
  \kappa(z_d)=v_d\kappa_2,
       \qquad v_d=N_{\bullet d}/8\in W .
 \tag{E17.7}
\]
To compare these framed divisors with tangential
paths, apply each boundary-unit character to the
framed Abel--Jacobi map. Regularization at the two
endpoints multiplies its value by
$(-1)^{m_d+m_O}$ for a divisor $\sum m_q q$.
This is the character of the torus tuple having
entry $+1$ at $d$ and $-1$ at every other non-$O$
cusp. Its coherent odd divisions are $H$-fixed.
Thus it contributes zero cocycle, proving (E17.7)
for the stated tangential paths as well.

Lift these abelian paths to full paths $p_d$ from
the tangent at $d$ to the tangent at $O$, composing
paths from right to left; the
map of path torsors is onto, since the quotient
$N\to N^{\rm ab}$ is onto. Put
$d_d=p_dt_dp_d^{-1}$ for positive cusp inertia.
The genuine conjugator $u_d(g)=p_dg(p_d)^{-1}$ has
abelian coordinate $v_d\kappa_2(g)$ and satisfies
$u_dg(d_d)u_d^{-1}=d_d^{\chi(g)}$.
The eleven peripheral classes together with a basis
of the compact summand form a basis of $N^{\rm ab}$;
their lifts form a free pro-$\ell$ basis by the
Frattini criterion. In the class-two coordinates of
E.15, put $L(e_d)=e_d\wedge v_d$ and $a=\kappa_2$.
The resulting actual $H$-action on the subgroup
$W\oplus\bigwedge^2W$ is
\[
 g(x,z)=(\chi_gx,\chi_g^2z+\chi_ga(g)L(x)).
 \tag{E17.8}
\]
Indeed the central term of
$u_d^{-1}d_d^{\chi_g}u_d$ is exactly the last term.
No action on a compact free generator is used here.

We give the averaging map used twice in the descent.
For an extension of $G_{\mathbb Q}$ by an abelian
$\mathbb Z_\ell$-module $M$ with a section $s_H$
over $H$, fix four left cosets $g_iH$ and lifts $e_i$.
For a lift $e$ of $g$ put $j=g\cdot i$,
$h_i=g_j^{-1}gg_i$, and
\[
 b_i(e)=e e_i s_H(h_i)^{-1}e_j^{-1}\in M,
 \quad B(e)=\sum_i b_i(e),
 \quad s_0(g)=(-B(e)/4)e .
 \tag{E17.9}
\]
Multiplying the displayed words gives
$B(ef)=B(e)+gB(f)$ and $B(m)=4m$.
Hence $s_0$ is a homomorphism independent of $e$.
It is continuous by using the continuous lifts
$e_i s_H(g_i^{-1}g)$ on the four clopen cosets.

First apply (E17.9) to the abelianized quartic
extension using the paths from the four $B_i$.
Its restriction to $H$ differs from the tangent
section at $O$ by $a(g)v$, with the positive sign
because $b_i(s_H(h))=-u_i(h)$ in abelian coordinates.
Both coordinate-unit rows annihilate (E17.6).
The leading values at $O$ are $\zeta_8$ and $1$,
whose odd Kummer cocycles on $H$ vanish. The toric
cocycle is therefore zero on $H$, and is zero
globally: a cocycle zero on this normal subgroup
has values in $\mathbb Z_\ell(1)^{H}=0$.
This constructs the asserted toric-zero abelian
origin $s_\ell$.

Over $H$, its literal class-two lift relative to
the original tangent has cocycle
\[
 c_v(g)=\bigl(a(g)v,\tfrac12a(g)^2L(v)\bigr).
 \tag{E17.10}
\]
Substitute (E17.8) and
$a(gh)=a(g)+\chi_ga(h)$ in the group law E15.1:
the extra term is
$\chi_ga(g)a(h)L(v)$, and the wedge of $v$ with
itself vanishes. This proves the cocycle identity.
Pull back the actual class-two extension along the
fixed $s_\ell$. The kernel is the abelian module
$\bigwedge^2N^{\rm ab}$, and (E17.10) gives its
$H$-section. Applying (E17.9) to this pullback
constructs $S_\ell$ over $G_{\mathbb Q}$.
This last averaging is in an abelian-kernel
extension, so its homomorphism proof applies
literally. The anchors are \S51.16, folios
733v--736v, scans 172--178, and task LM7-E.23(iv).
\end{proof}

\subsection{The cubic class-three map and its remaining arithmetic input}

\begin{proposition}[E.18 --- Proved: the exact class-three criterion and every nonempty fibre]
\label{E-class-three-criterion}
Let $N=\pi_1(U_{3,\overline{\mathbb Q}})^{(3)}$, with
the genuine action defined by the rational unit tangent
at $O=[1:-1:0]$. Put $N_i=N/\gamma_{i+1}N$,
$V=N^{\rm ab}$, $A=\gamma_3N/\gamma_4N=\operatorname{Lie}_3V$,
and $K=\mathbb Q(\zeta_3)$. The rank of $A$ is $330$.
Fix a continuous ordered-Hall set section $h:N_2\to N_3$
carrying one to one. For an actual continuous cocycle
$a:G_{\mathbb Q}\to N_2$ and every prime
$\mathfrak q\nmid3$ of $K$, choose a prolongation,
positive tame generator $t_{\mathfrak q}$ and arithmetic
Frobenius $F_{\mathfrak q}$. Write $q=N\mathfrak q$ and
\[
 \begin{aligned}
 n_{\mathfrak q}&=h(a(t_{\mathfrak q})),&
 x_{\mathfrak q}&=h(a(F_{\mathfrak q})),\\
 \epsilon_{\mathfrak q}
 &=x_{\mathfrak q}\rho_3(F_{\mathfrak q})(n_{\mathfrak q})
       x_{\mathfrak q}^{-1}n_{\mathfrak q}^{-q}\in A,&
 z_{\mathfrak q}
 &=-(\Phi_{\mathfrak q}-q)^{-1}\epsilon_{\mathfrak q}
       \in A\otimes\mathbb Q_3 .
 \end{aligned}
 \tag{E18.1}
\]
Here $\rho_3$ is the actual action on $N_3$ and
$\Phi_{\mathfrak q}$ its induced action on $A$.
The fibre over $[a]$ of
\[
 H^1_{\rm cts}(G_{\mathbb Q},N_3)
       \longrightarrow H^1_{\rm cts}(G_{\mathbb Q},N_2)
 \tag{E18.2}
\]
is nonempty precisely when every $z_{\mathfrak q}$ is
integral and $z_{\mathfrak q}\to0$ in the uniform
$3$-adic congruence topology. Each nonempty fibre is
a torsor under the explicitly evaluated group
\[
 \mathbb Z_3^{166}\oplus(\mathbb Z/3)^7
   \oplus(\mathbb Z/9)^{56}\oplus(\mathbb Z/27)^{25}
   \oplus(\mathbb Z/81)^3 .
 \tag{E18.3}
\]
For every class unramified outside $3$, this fibre
is nonempty. At the $K$-prime $(2)$ the integral
local obstruction module is
\[
 \coker(\Phi_{(2)}-4)=
 (\mathbb Z/3)^{70}\oplus(\mathbb Z/9)^{24}
 \oplus(\mathbb Z/27)^{76}\oplus(\mathbb Z/81)^{28}.
 \tag{E18.4}
\]
The test (E18.1) retains the higher action and the
actual Frobenius cocycle values as data. Its proof
and the two computed coefficient groups verify LM7
E.41 and E.44 in the degree used here. E.21 records
the unevaluated arithmetic data for unrestricted
inputs rather than counting this test as their
complete evaluation.
\end{proposition}
\begin{proof}
The ordered commutator basis of a free pro-three
group gives $\gamma_3N/\gamma_4N=\operatorname{Lie}_3V$.
In the convention $[[x_i,x_j],x_k]$ with $i>j$ and
$k\ge j$, it has $(10^3-10)/3=330$ basis elements.
Ordered products with $3$-adic exponents supply the
continuous Hall set section; no BCH division by
three is required.

The set lift $\widehat a=h\circ a$ has central defect
\[
 \omega(g,h)=\widehat a(g)g(\widehat a(h))
                    \widehat a(gh)^{-1}\in A.
 \tag{E18.5}
\]
Associativity proves the two-cocycle identity.
A correcting continuous cochain $b$ gives a lift
exactly when $db=-\omega$. At a prime away from
three the genuine action is unramified. Wild
inertia has trivial image in the pro-three
geometric group. The lower inertia cocycle is a
homomorphism, and we lift it by
$t^u\mapsto n^u$ for $u\in\mathbb Z_3$.
The tame relation gives precisely
$\epsilon_{\mathfrak q}$ in (E18.1).
Changing $n$ to $zn$ changes this error by
$(\Phi_{\mathfrak q}-q)z$; changing $x$ centrally
has no effect. The rational weights of $A$ have
Frobenius absolute values $q^{j/2}$ for
$3\le j\le6$. None equals $q$, proving existence
and uniqueness of the inverse used in (E18.1).

Here is the continuous global comparison, including
its image rather than its closure. Define
\[
 \mathcal R(A)=\{(v_{\mathfrak q})\in
          \prod_{\mathfrak q\nmid3}A:
                   v_{\mathfrak q}\longrightarrow0\},
 \qquad d(v)_{\mathfrak q}=(q^{-1}\Phi_{\mathfrak q}-1)v_{\mathfrak q}.
\]
We have the actual isomorphism and representative
\[
 H^2_{\rm cts}(G_K,A)=\mathcal R(A)/d\mathcal R(A),
 \qquad [\omega]\longmapsto
       [(q^{-1}\epsilon_{\mathfrak q})_{\mathfrak q}].
 \tag{E18.6}
\]
To prove it, let $S=\operatorname{Spec}\mathcal O_K[1/3]$
and $A_m=A/3^mA$. The proof of E.16 gives
$H^2(G_{K,\{3\}},A_m)=0$; the action has pro-three
image and a filtration with trivial simple factors.
The number-field cohomological-dimension theorem,
the finiteness theorem, and the scheme comparison
give $H^3(S,A_m)=0$ and finite $H^1(S,A_m)$.
We use exactly Milne I.4.10(c), I.4.15 and II.2.9,
with three inverted and no real place
\cite{E_density_MilneADT}.

For a finite set $J$ of primes away from three,
localization and purity give the triangle
\[
 R\Gamma(S,A_m)\longrightarrow
 R\Gamma(S-J,A_m)\longrightarrow
 \bigoplus_{\mathfrak q\in J}
 R\Gamma(k_{\mathfrak q},A_m(-1))[-1]
 \longrightarrow R\Gamma(S,A_m)[1].
 \tag{E18.7}
\]
The local purity term is also obtained directly:
over a strict henselization wild inertia has zero
three-primary cohomology and tame inertia is
$\mathbb Z_3(1)$. Take the direct limit over finite
$J$, then the derived inverse limit in $m$.
The residue complexes are the compatible two-term
complexes $[A_m\xrightarrow{q^{-1}\Phi-1}A_m]$.
Both direct-sum term systems have surjective
transition maps. Their inverse limit is the literal
complex $[\mathcal R(A)\xrightarrow d\mathcal R(A)]$.
Finite $H^1(S,A_m)$ has vanishing first derived
limit, and the two higher unramified terms vanish.
Jannsen's Proposition 1.6 and Theorem 2.2 identify
this calculation with continuous-cochain cohomology
\cite{Jannsen1988}. This proves the first equality
in (E18.6).

The support map measures failure of the chosen
inertia splitting to be Frobenius invariant.
Evaluating its transport on $t^{q^{-1}}$ gives
$q^{-1}\epsilon_{\mathfrak q}$, which proves the
second formula and its Tate twist. Its family is
restricted: at every fixed finite characteristic
quotient the cocycle graph has finite image and
is ramified at finitely many primes. Thus all its
inertia values approach one uniformly outside
finite sets. Continuity of $h$ and of powers
$n^u$ on the compact spaces in question gives
the same property for the local null-homotopies
used in (E18.7). Hence their errors are zero
modulo $3^mA$ outside a finite set for every $m$.

Equation (E18.6), with the rational invertibility
already proved, makes its vanishing equivalent
to the integral restricted corrections of the
statement. Restriction from $G_{\mathbb Q}$ to
$G_K$ is injective on this $H^2$: restriction
followed by corestriction is multiplication by
two, an automorphism. Thus the same criterion
gives a continuous global lift. Changing Hall
lifts changes the support sequence by an actual
element of $d\mathcal R(A)$. Changing tame
generators rescales by units, changing Frobenius
by an inertia power preserves its error, and
changing prolongations transports by integral
Galois automorphisms followed by such a boundary.
The criterion therefore depends only on $[a]$.

Differences of two lifts are $A$-valued cocycles.
At a good prime their inertia values satisfy
$(\Phi-q)v=0$, so they vanish. Every such
cocycle is therefore unramified outside three.
There is also no lower geometric stabilizer
to divide out: the lowest nonzero graded
coordinate of a fixed element of the
section-twisted $N_2$ would be fixed by the
arithmetic action. At any good Frobenius all
its eigenvalues have absolute value greater
than one; inner twisting acts trivially on
this lowest coordinate. The contradiction
proves triviality of the stabilizer. The fibre
is exactly a torsor under $H^1(G_{\mathbb Q,\{3,\infty\}},A)$.
E.16's embedding-problem argument lifts every
unramified lower section through this central
step as well.

We supply the integral coefficient evaluation.
The eight nonorigin cusps are
$A_i=[0:-\zeta^i:1]$, $B_i=[-\zeta^i:0:1]$,
and $C_1,C_2=[1:-\zeta^i:0]$.
In the order $A_0,A_1,A_2,B_0,B_1,B_2,C_1,C_2$,
their elliptic three-torsion coordinates are
\[
 (u_d,v_d)=(1,0),(2,2),(0,1),(2,0),
                 (1,1),(0,2),(2,1),(1,2).
 \tag{E18.8}
\]
Tangent-line divisors make all nine cusps three-torsion;
their collinear triples give this addition table.
On $Y^2+Y=X^3-7$ the basis points are
$(3,-5),(3\zeta,-5)$, whose tangent-line ratio gives
the pairing $-9(\zeta-1)/(9(\zeta^2-1))=\zeta$.
For the framed generalized-Jacobian extension,
translation on the trivialization of the pullback
line bundle is the inverse of this pairing.
Thus, writing $U_d=(v_d,-u_d)$, the integral
Tate map of the unique $j_0$ with $pj_0=[3]$ is
$j_{0*}=\binom{3I_2}{U}$. Equivariance forces
the following actual degree-one matrices:
\[
 \begin{gathered}
 R=\begin{pmatrix}0&-1\\1&-1\end{pmatrix},\quad
 C_E=\begin{pmatrix}1&-1\\0&-1\end{pmatrix},\\
 A(\psi,q)=
 \begin{pmatrix}A_E&0\\U(A_E-qI)/3&qI_8\end{pmatrix},
       \quad A_E=aI+bR,\quad\psi=a+b\zeta,\\
 C_V=\begin{pmatrix}C_E&0\\(UC_E-C_TU)/3&C_T\end{pmatrix}.
 \end{gathered}
 \tag{E18.9}
\]
Here $C_T$ is minus the permutation interchanging
subscripts one and two. All displayed divided
blocks are integral. The elliptic Frobenius
scalars at $(2)$ and a specified prime over seven
are $-2$ and $1+3\zeta$: their norm, point count,
and primary congruence determine them. In
$1+3\mathbb Z_3[\zeta]$, the logarithms divided
by three of $1+3\zeta$ and $1+3\zeta^2$ reduce
to a basis modulo three. They therefore give
the full rank-two CM image. Since the unramified
pro-three source is free of rank two, a Frobenius
lift $\tau$ at the chosen prime and
$\tau'=c\tau c$ form a free basis by the
Frattini criterion.

An $A$-cocycle normalized to vanish at $c$ is
freely determined by its value at $\tau$.
Its value at $\tau'$ is $J$ times that value,
where $J=\operatorname{Lie}_3(C_V)$.
The coboundaries preserving this normalization
are $(T-1)A^c$, where
$T=\operatorname{Lie}_3(A(1+3\zeta,7))$.
Since two is a unit, $A^c=(1+J)A$. Hence the
actual fibre group is
\[
 \coker\bigl((T-1)(1+J):A\to A\bigr).
 \tag{E18.10}
\]
The program \texttt{class\_three\_fibre.py}
constructs both integral matrices in the Hall
basis. It verifies $J^2=1$, $\tr J=-2$, and
the Smith valuations
$0^{73},1^7,2^{56},3^{25},4^3$.
These are all nonzero elementary divisors:
$1+J$ has exact rational rank $164$, and the
calculation finds $164$ pivots at precision
$3^{10}$. Its remaining free rank is $166$.
This proves (E18.3), including its torsion.

For (E18.4), the independent program
\texttt{class\_three\_matrix.py} applies the
same Hall functor to $A(-2,4)$. It finds all
$330$ pivots of $\Phi-4$ at precision $3^9$,
with valuations
$0^{132},1^{70},2^{24},3^{76},4^{28}$.
Each pivot is strictly below the working
precision, so integral elementary operations
lift with the same valuations. As a separate
check, the rational eigenvalues have
multiplicities
$64^{168},(-32)^{128},16^{32},(-8)^2$;
their determinant difference from four has
valuation $168+256+32+2=458$, equal to the
sum of the computed valuations. This proves
the full cokernel, of order $3^{458}$.
All source maps, actions, and choices have
therefore been retained. The manuscript
anchors are \S51.16, folios 733v--736v,
scans 172--178, and LM7-E.23(i).
\end{proof}

\begin{proposition}[E.19 --- Proved: a linear reduction of the entire zero-abelian slice]
\label{E-zero-slice}
Retain E.18 and put $M=\bigwedge^2V$ and
$L=\gamma_2N/\gamma_4N$. The fibre of
$H^1(G_{\mathbb Q},N_2)\to H^1(G_{\mathbb Q},V)$
over zero is exactly $H^1(G_{\mathbb Q},M)$.
Its class-three obstruction is the linear connecting
homomorphism for the actual $G_{\mathbb Q}$-module
sequence
\[
 0\longrightarrow A\longrightarrow L
       \longrightarrow M\longrightarrow0 .
 \tag{E19.1}
\]
This whole obstruction is determined by the action
on $N_2$. An action on $N_3$ need not be supplied
separately for this slice.

More precisely let
$w=(3f_1+U_1)\wedge(3f_2+U_2)\in M$,
in the integral coordinates (E18.9).
It is primitive and spans a $G_{\mathbb Q}$-submodule
$\mathbb Z_3(1)$. Every $[z]\in H^1(G_{\mathbb Q},M)$
has a unique decomposition
\[
 [z]=[z_0]+\beta w,\qquad
 [z_0]\in H^1(G_{\mathbb Q,\{3,\infty\}},M),
 \qquad \beta=\sum_{p\ne3}a_p\kappa_3(p),
 \quad a_p\in\mathbb Z_3,\quad a_p\to0.
 \tag{E19.2}
\]
Choose a lift $\omega\in L$ of $w$ with
$c\omega=-\omega$ for complex conjugation $c$.
The actual continuous cocycle
\[
 \eta(g)=\chi(g)^{-1}\bigl(g\omega-\chi(g)\omega\bigr)
       \in A(-1)
 \tag{E19.3}
\]
is unramified outside three. The obstruction of
(E19.2) is exactly
\[
                  [\eta\cup\beta]\in
                        H^2_{\rm cts}(G_{\mathbb Q},A).
 \tag{E19.4}
\]
Over $K$, its unique rational tame corrections are
\[
 -a_p\,(q^{-1}\Phi_{\mathfrak q}-1)^{-1}
                          \eta(F_{\mathfrak q}),
       \qquad \mathfrak q\mid p,\quad p\ne3.
 \tag{E19.5}
\]
Thus the exact image in this slice is the kernel
of the linear map (E19.4), equivalently the integral
restricted locus of (E19.5); its nonempty fibres
are (E18.3).

The remaining action class (E19.3) is a single
element of the explicitly computed group
\[
 H^1(G_{\mathbb Q,\{3,\infty\}},A(-1))
 =\mathbb Z_3^{164}\oplus(\mathbb Z/3)^{59}
       \oplus(\mathbb Z/9)^{27}\oplus(\mathbb Z/27)^7.
 \tag{E19.6}
\]
Its normalized cocycle is determined by the one
vector $\eta(\tau)\in A(-1)$, with $\tau$ as in
E.18. We have evaluated the coefficient group
and the exact reduction maps, not this vector.
The unresolved evaluation is retained in E.21.
\end{proposition}
\begin{proof}
An abelianly trivial class in $H^1(N_2)$ can be
conjugated by a lift of its abelian coboundary to
a cocycle in $M$. Its further ambiguity from
$V^{G_{\mathbb Q}}$ vanishes by good Frobenius
weights, exactly as in E.18. This proves the
first identification. The group $L$ is abelian
because $[\gamma_2N,\gamma_2N]\subset\gamma_4N$.
Its two graded pieces are free $\mathbb Z_3$-modules,
and ordered Hall products make (E19.1) a sequence
of finite free modules with continuous action.
A linear Hall section $\iota:M\to L$ gives
\[
 g\iota(m)=\iota(B_gm)+J_gm,
       \qquad B_g=\bigwedge^2 A_g,
       \quad J_g:M\to A.
\]
For an $M$-cocycle $z$, its lift has defect
$J_gz(h)$, the connecting cocycle of (E19.1).
This proves linearity. The action on every
commutator $[x_i,x_j]$ modulo $\gamma_4N$
depends only on the images of $x_i,x_j$ modulo
$\gamma_3N$: changing either by a degree-three
element changes this commutator in degree four.
It follows on the Hall generators that the
entire action on $L$ is determined by the
given action on $N_2$.

The coefficient of $w$ on
$e_{A_0}\wedge e_{A_2}$ is one, so $w$ is primitive.
The map $j_{0*}$ in E.18 is equivariant, and its
exterior square has determinant character $\chi$.
At each good prime the only rational eigenline
of $M$ with Frobenius eigenvalue $p$ is this
compact determinant line: all other eigenvalues
have absolute values $p^{3/2}$ or $p^2$.
Consequently the integral inertia value of an
$M$-cocycle is $a_pw$. These coefficients tend
to zero by continuity. The uniformly convergent
sum $\sum_{p\ne3}a_p\kappa_3(p)w$ is an actual
continuous cocycle with precisely these residues.
Its subtraction kills all inertia away from
three, giving (E19.2). The residues and the
remaining unramified class also prove uniqueness.

Average an arbitrary lift of $w$ as
$\omega=(\omega_0-c\omega_0)/2$.
This still lifts $w$ and satisfies the desired
sign. Formula (E19.3) then has $\eta(c)=0$.
Applying $gh$ to $\omega$ shows
$\eta(gh)=\eta(g)+\chi(g)^{-1}g\eta(h)$,
which is exactly the cocycle law in $A(-1)$.
The action on $L$ is unramified away from
three, proving the same for $\eta$.
Changing $\omega$ adds a coboundary, so its
class is the pullback of (E19.1) along
$\mathbb Z_3(1)w\hookrightarrow M$.

The obstruction of $z_0$ vanishes by
$H^2(G_{\mathbb Q,\{3,\infty\}},A)=0$,
proved in E.16. Lift the other term by
$g\mapsto\beta(g)\omega$. Its defect is
$\chi(g)\beta(h)\eta(g)$, which is the
ordered cup $\eta\cup\beta$ under
$A(-1)\otimes\mathbb Z_3(1)\to A$.
This proves (E19.4) with its sign. Locally,
both the inertia and Frobenius lift lie in
the abelian group $L$. Their conjugation
term in (E18.1) therefore disappears, and
the error is
$a_p(g\omega-q\omega)=qa_p\eta(F)$.
Equation (E18.1) gives (E19.5).

Finally the same free source basis as in
E.18 evaluates the group containing $\eta$.
The matrices on $A(-1)$ are $T/7$ and
$-J$. The normalized cocycle is determined
freely by its value at $\tau$, and its
coboundaries form
$(T-7)(1-J)A$, since seven is a unit.
The independent computation
\texttt{class\_three\_zero\_slice.py}
constructs this matrix from E18.9. It finds
$166$ nonzero pivots with Smith valuations
$0^{73},1^{59},2^{27},3^7$ at precision
$3^{10}$. These are all its elementary
divisors because $1-J$ has exact rank
$(330+2)/2=166$. The cokernel has the
$164$ free coordinates in (E19.6).
This calculation specifies the finite
arithmetic datum still to evaluate; it
does not infer its value from the matrix
of the graded action. The manuscript
anchors are \S51.16, folios 733v--736v,
scans 172--178, task LM7-E.23(i).
\end{proof}

\begin{proposition}[E.20 --- Proved: a fully checked class-two lift obstructed at class three]
\label{E-finite-next-obstruction}
Let $P=[0:-1:1]$, $Q=[-1:0:1]$ on the plus Fermat
cubic, and let $\alpha_P\in H^1(G_{\mathbb Q},V)$ be
the class of the rational tangent at $P$ with value
two in the parameter $-x/z$. For positive cusp
inertia $e_Q\in V$, set
\[
           \alpha=\alpha_P+27\kappa_3(19)e_Q.
 \tag{E20.1}
\]
This class lifts to the full intrinsic class-two
quotient of $N$. Every such lift fails to extend
to $N/\gamma_4N$, $N/N''$, or $N/\gamma_3I_Q$,
where $I_Q$ is the closed normal subgroup generated
by the marked cusp $Q$. The failure is detected
at $19$ by an actual arithmetic quotient of $N$
of order $3^{30}$ and nilpotency class three.
This independently verifies the exact example
of LM7 E.38 and its transfer maps.
\end{proposition}
\begin{proof}
We first check existence at full intrinsic class
two. Use $K,V,M,w$ and the integral lattice of
E.18--E.19, and put $M_0=\ker(M\to\bigwedge^2T_3E)$.
Since $w$ has coefficient one at
$e_{A_0}\wedge e_{A_2}$, deleting that column gives
a basis of $M/\mathbb Z_3w$. Deleting the row
$f_1\wedge f_2$ gives a basis of $M_0$.
At a prime of norm $q$ the map
\[
 d_{\mathfrak q}:M/\mathbb Z_3w\longrightarrow M_0,
       \qquad [m]\longmapsto(q^{-1}\bigwedge^2A_{\mathfrak q}-1)m
 \tag{E20.2}
\]
is rationally invertible, by the weight computation
in E.19. The same localization proof as E.18,
now in degree two and retaining its one-dimensional
rational kernel, identifies the obstruction with
the restricted support cokernel of these maps.

The point-tangent class $\alpha_P$ lifts to the
full group. Its framed cusp-path values have
support only at three, and its tangent value
adds the prime two. This follows directly by
evaluating the ratios of the cubic tangent lines;
their nonzero values are products of roots of
unity and $1-\zeta_3$. At the prime of $K$ over
two, $19\in1+2\mathcal O_{K,2}$ has coherent
third-power roots because cubing is an automorphism
of this pro-two unit group. Thus the added
Kummer class is zero locally, and the obstruction
there is unchanged from $\alpha_P$.

At a prime over $19$, normalize
$\kappa_3(19)(t)=1$. Let $a_P=\alpha_P(F)$, and
let $k_Q$ be the abelian path conjugator from the
unit tangent at $Q$ to the fixed origin. Expanding
the class-two action on its inertia loop and the
tame relation gives the support numerator
\[
 D=(a_P+27\kappa_3(19)(F)e_Q)\wedge27e_Q
                   -27k_Q\wedge e_Q
       =27(a_P-k_Q)\wedge e_Q\in27M_0.
 \tag{E20.3}
\]
The two elliptic Frobenius scalars over $19$ are
$-2+3\zeta$ and $-5-3\zeta$. Substitution in
E18.9 and exterior squaring construct the two
actual $44$-by-$44$ matrices (E20.2).
The independently executed program
\texttt{class\_three\_counterexample.py} finds
Smith valuations $0^8,1^6,2^{18},3^{12}$ for each.
All $44$ pivots are found below its precision,
and their valuation sum is $78$. In particular
$27M_0\subset d_{\mathfrak q}(M/\mathbb Z_3w)$,
so (E20.3) has an integral correction at both
primes. Every other good-prime support term is
zero. These corrections have finite support,
and the continuous localization quotient
therefore gives a global class-two lift.

For the higher obstruction use the unit tangent
at $P$. The power-cover coordinate $t=s^3$,
$s=-x/z$, and the branch
$y/z=-(1-s^3)^{1/3}$ with rational leading
coefficient give a genuine equivariant inclusion
\[
 N=\ker(F\to(\mathbb Z/3)^2)\hookrightarrow
       F=\langle x,y\rangle^{\rm pro\text{-}3}.
 \tag{E20.4}
\]
A path to the previous tangent changes the
arithmetic origin by its actual tangential
cocycle and preserves the defined abelian
class (E20.1). Put $I_y=\ker(F\to\mathbb Z_3)$,
where $x\mapsto1$, $y\mapsto0$. With $T=1+U$,
the exact quotient map is
\[
 \begin{gathered}
 F/[I_y,I_y]=R=\mathbb Z_3\ltimes\mathbb Z_3[[U]],\\
 (a,v)(b,w)=(a+b,v+T^a w),\qquad
 x\mapsto(1,0),\quad y\mapsto(0,1),\\
 N\twoheadrightarrow H=3\mathbb Z_3\ltimes J,
            \qquad J=(3,U)\mathbb Z_3[[U]].
 \end{gathered}
 \tag{E20.5}
\]
To verify the quotient, the conjugates of $y$
generate $I_y^{\rm ab}$ over the completed
group ring $\mathbb Z_3[[\mathbb Z_3]]$.
The universal free-group map to the displayed
semidirect product is inverse on this module
generator. It proves equality of the modules,
the group law, and the topology. On
$F^{\rm ab}$ the map is the identity, with
second coordinate the augmentation $v(1)$;
this proves the stated image of $N$.

The normalized tripod action is
$x\mapsto x^\chi$, $y\mapsto f_g^{-1}y^\chi f_g$
with $f_g\in F'$. In $R$, this conjugator lies
in the abelian module and commutes with $y$.
Consequently the actual descended action is
\[
           g(a,v(T))=(\chi(g)a,\chi(g)v(T^{\chi(g)})).
 \tag{E20.6}
\]
The group $R$ is metabelian. Every loop above
the base puncture one maps into its abelian
normal module. Hence the same map from $N$
factors through $N/N''$ and $N/\gamma_3I_Q$,
with all arithmetic actions retained.

At $19$, choose $\zeta_9$ reducing to
$4\in\mathbb F_{19}$; then
$2^{(19-1)/9}=4$ gives
$\kappa_3(2)(F)\equiv1\pmod9$.
The first coordinate of any section above
(E20.1) is $3\kappa_3(2)$, and its tame inertia
has augmentation $3\cdot27=81$.
Conjugating by $x^{3c}$ normalizes its Frobenius
first coordinate to three: the change is
$3(1-19)c$, and
$(\kappa_3(2)(F)-1)/18$ is integral.
The inertia value is therefore $(0,v(T))$
with $v(1)=81$, and the tame relation becomes
\[
                         T^3v(T^{19})=v(T).
 \tag{E20.7}
\]
The Frobenius module coordinate cancels, so this
equation covers every possible choice of lift.
Over $\mathbb Q_3[[U]]$, its unique solution is
$v(T)=81T^{-1/6}$. Indeed putting $T=e^z$ and
multiplying by $e^{z/6}$ makes every positive
coefficient invariant under $z\mapsto19z$,
so it vanishes. Its fourth coefficient is
$1729/384$, of three-adic valuation $-1$.

For a finite certificate, reduce the module
modulo $(729,U^5)$ and the first coordinate
modulo $2187$. The coefficients of
$(1+U)^{2187}-1$ through degree four are
divisible by $729$, so (E20.5)--(E20.6)
descend to a finite arithmetic quotient.
Its subgroup $H_6$ has order $3^{35}$.
Writing $v=81+\sum_{i=1}^4v_iU^i$ gives
$Kv=b\pmod{729}$, where
\[
 K=\begin{pmatrix}
 18&0&0&0\\228&360&0&0\\
 1539&7581&6858&0\\7315&86640&205770&130320
 \end{pmatrix},\qquad
 b=\begin{pmatrix}-243\\-243\\-81\\0\end{pmatrix}.
 \tag{E20.8}
\]
The row $r=(1,660,558,162)$ satisfies
$rK=0$ and $rb=486$ modulo $729$.
These identities are reproduced by the program
directly from the polynomial substitution.

Let $\delta=T^3-1$. Commutators in $H_6$ give
$\gamma_{c+1}H_6=\delta^cJ$. The subgroup
$\gamma_4H_6$ is generated by
\[
       3\delta^3=81U^3+243U^4,
               \qquad U\delta^3=27U^4.
 \tag{E20.9}
\]
It has order $3^5$: the first generator has
order nine, the second order twenty-seven,
and their cyclic groups meet trivially by
their $U^3$ coordinates. The program also
enumerates all $243$ elements. The row $r$
annihilates both columns in (E20.9), so its
nonzero pairing with $b$ persists in the
quotient $H_6/\gamma_4H_6$ of order $3^{30}$.
This quotient has class exactly three,
because $\delta^2U=9U^3+18U^4$ is not in
(E20.9). The arithmetic surjection
\[
             N\twoheadrightarrow H
                 \twoheadrightarrow H_6/\gamma_4H_6
 \tag{E20.10}
\]
factors through $N/\gamma_4N$ by functoriality
of the lower central series and through the
other two specified quotients by (E20.5).
Every proposed section would yield the
excluded local solution of (E20.8).
This proves all three failures. The anchors
are \S51.16, folios 733v--736v, scans
172--178, and LM7-E.23(i)--(ii).
\end{proof}

\begin{proposition}[E.21 --- Open: the unrestricted intrinsic tower image]
\label{E-remaining-tower}
The following requested classifications remain
unfinished in this return:
\begin{enumerate}[label=(\roman*)]
\item The evaluated image of (E18.2) for every
ramified class-two input on the cubic. Its nonempty
fibres are (E18.3), and its exact continuous test
is (E18.1). The action vector $\eta(\tau)$
of E.19 is partially evaluated in E.35--E.37:
its two rational Kummer coordinates are known,
its elliptic projection is an explicit Ross
cup edge, and one primitive integral row has
value six. The numerical Ross value and
integral torsion remain.
\item The image and all fibres of the maps from
the intrinsic cubic metabelian and marked-depth
quotients to the inherited quotient $H$ of
(E20.5), for every first class.
\item The complete metabelian section sets for
the other coefficient primes and for finite
quotients involving several primes, including
all odd Fermat levels. This includes the first
level in the full profinite nilpotent tower
capable of selecting precisely the cuspidal
classes at an odd level.
\end{enumerate}
Task LM7-E.23(iv), the quartic local Kummer-order
law and additive defect at $3$ and $5$, is
evaluated by E.13--E.15 with the actual origin
verified in E.17. That conclusion leaves
(i)--(iii) on the page.
\end{proposition}
\begin{proof}[State reached and the attempted steps]
We checked the supplied LM7 E.37, E.38, E.41
and E.44 proofs where used. Their later programs
were absent from the supplied expanded state;
the matrices and finite certificate reported
here were constructed and executed independently.
The earlier toric and local class-two criteria
are retained through the exact quotient maps
in E.1, E.15 and E.20.

For (i), replacing $\rho_3(F)$ in (E18.1)
by $\operatorname{Lie}_{\le3}(A_F)$ has no
justification: the displayed support error
uses the images of actual generators, whereas
the graded representation records only their
leading coordinates. We retained that term
and computed every nonempty fibre independently.
E.20 evaluates a ramified input with a finite
negative answer; E.16 gives an entire positive
unramified family. E.19 goes further on the
whole zero-abelian slice: the obstruction
factors through the class of (E19.3), so only
the action on degree-two commutators and one
normalized global vector remain. The original next
specific computation was the image of the
primitive compact determinant $w$ under the
actual Frobenius at the chosen prime over
seven, through $\gamma_2N/\gamma_4N$.
The peripheral path values of LM7 E.37
do not supply the images of both compact
free generators needed for this vector.
Passing to their linear Tate matrices gives
the coefficient group (E19.6), but leaves
the vector undetermined.

For (ii), the kernel of the exact map
$N\to H$ is $[I_y,I_y]\subset N$.
Thus the two kernel groups to be restored are
$[I_y,I_y]/N''$ and
$[I_y,I_y]/\gamma_3I_Q$, respectively.
A section of the target defines a pulled-back
extension by that specified kernel; its
splitting is the precise next question.
The local calculations of radical-class
towers in the supplied E.111, E.114 and
E.116 compute target sections. They do
not provide such splittings, and E.20
shows that passage back can fail already
in a finite class-three quotient.
The next approach is to compute the first
central graded kernel in these pullbacks
for the target families, retaining their
Frobenius correction words in E18.1.

For (iii), the even plus levels are settled
at their integral geometric abelian
two-primary quotient by E.5. For every
$n=3^r$ and $n=5^r$, E.16 proves that no
intrinsic nilpotent quotient at that
coefficient prime selects only cusps:
the constructed full pro-prime sections
have infinite-order compactified images.
The map from the full geometric profinite
group to a product of its maximal pro-prime
quotients kills elements detected by finite
groups involving several primes. A lift
from that product to the full metabelian
quotient must therefore solve its actual
kernel extension; none of the separate
coefficient calculations supplies it.
We have neither classified this lift
image nor proved a universal first
nilpotent level for odd $n$.
E.31 below makes one further comparison explicit:
the characteristic Haar summand embeds the entire
three-primary derived module into the full
profinite derived module. It supplies an actual
$\mathbb Z_3$-family of zero-abelian noncuspidal
full metabelian sections at the cubic level.
The complete mixed-prime image and its other
fibres remain unevaluated here.
E.32 also assembles the E.16 section with
the base cusp sections at all other primes
in the full pronilpotent product. It proves
that no full nilpotent level works for
$n=3^r$ or $5^r$. E.33 constructs the larger
full metabelian zero-abelian families for
every $n=3^r$, with their exact ranks;
E.34 proves the same full-metabelian family
statement for every $n=5^r$. E.35--E.36
identify the remaining elliptic class and
evaluate both rational Kummer coordinates.
E.37 gives the zero-abelian class
$\kappa_3(19)w$ a finite coefficient
obstruction at nineteen. E.38 constructs an
intrinsic zero-abelian class with lifts at
every finite coefficient level and no
continuous class-three lift. E.39 computes
the exact finite kernel lost by the compact
and unit cohomology comparisons. The actual
integral path class still requires correlated
division-point and unit-root data, so arbitrary
cochain representatives cannot fill that step.
These results retain the remaining complete image and
fibres as the open classification task.

For (iv), the exact density calculation
uses primitive integral Tate cocycles,
proves their full affine image, counts
finite Cartan and non-Cartan strata,
and controls the determinant-zero
boundary by measure tending to zero.
E.15 proves the polarization in the
genuine arithmetic action and its
restricted mixed correction. Thus the
local-order law, its scalar shifts,
and the failure-of-additivity pairing
are theorems; they do not evaluate
the different remaining input of (i).
The manuscript anchors for all four
tasks are \S51.16, folios 733v--736v,
scans 172--178.
\end{proof}

\begin{proposition}[E.22 --- Proved: the precise comparison with unipotent Chabauty]
\label{E-Kim-comparison}
Fix a rational point or rational tangent on a
hyperbolic curve $Y/\mathbb Q$, a coefficient
prime $\ell$, and a nilpotency class $c$.
The actual continuous coefficient maps are
\[
 \pi_1(Y_{\overline{\mathbb Q}})
 \longrightarrow N^{(\ell)}
 \longrightarrow N^{(\ell)}/\gamma_{c+1}
 \longrightarrow U_{\ell,c}(\mathbb Q_\ell),
 \tag{E22.1}
\]
where $U_{\ell,c}$ is the class-$c$ unipotent
completion. They induce maps of section
classes, compatible with every restriction
to $G_{\mathbb Q_p}$ and with the point
Kummer maps. For a regular model
$\mathcal Y/\mathbb Z_S$, $\ell\notin S$,
the image of the $S$-integral finite-descent
set at $\ell$ is consequently contained
in the refined Chabauty--Kim locus for
every unipotent quotient.

Under the coefficient map $A\to A\otimes\mathbb Q_3$
of E.18, every finite-support obstruction
vanishes. Unrestricted support retains
the boundedness and convergence test.
In particular the finite integral local
obstruction modules (E18.4) and (E20.8)
map to zero under this specified
rationalization. The local point-image
conditions used by Chabauty--Kim remain
additional conditions on the target
cohomology classes, defined by the
commuting maps just stated.
\end{proposition}
\begin{proof}
For a free geometric pro-$\ell$ group, map
its free generators to $1+X_i$ in the
completed noncommutative power-series
algebra. Modulo words of degree $c+1$,
the Magnus map sends group elements to
group-like units; finite logarithm and
exponential polynomials over $\mathbb Q_\ell$
give the corresponding unipotent group.
Ordered Hall products show directly
that its map on each graded piece is
the usual lattice inclusion into the
free rational Lie algebra. This
constructs (E22.1); for a proper curve,
quotient by the closed surface relation
and its unipotent image. The universal
property extends each arithmetic
automorphism, so the maps are equivariant.
Applying them to a continuous cocycle
commutes literally with restricting
its domain. Applying them to a path
torsor gives the corresponding
unipotent path torsor. These two
constructions prove the asserted
commuting Kummer maps.

An $S$-integral adelic point surviving
finite descent is the localization of
a Selmer section by the Harari--Stix
comparison used in D.1. Push this
section through (E22.1). At every finite
place its restriction lies in the
unipotent Kummer image of the chosen
local point, and outside $S$ in that
of integral points. Thus its global
class satisfies the defining refined
Selmer local conditions. Its value
at $\ell$ is in the scheme-theoretic
localization image, proving the
inclusion. This is the exact map
in Betts--Kumpitsch--L\"udtke,
Lemma 2.19, whose statement and proof
we checked~\cite{BettsKumpitschLudtke2026}.
It applies in particular to the odd
Fermat curves with the rational
tangent used above.

For the integral-to-rational assertion,
every $\Phi_{\mathfrak q}-q$ in E.18
is rationally invertible. A support
family at finitely many primes has
only finitely many rational inverse
vectors, so one power of three makes
all of them integral. The continuous
obstruction is therefore torsion and
maps to zero after tensoring with
$\mathbb Q_3$. More explicitly, a
continuous cochain from a compact
profinite domain to $A\otimes\mathbb Q_3$
has bounded image. Hence its cochain
complex is the filtered union of
the complexes with values in $3^{-n}A$,
and its cohomology is the localization
of the integral continuous cohomology.
For an infinite support sequence,
the same map requires a common
denominator and convergence to zero;
these requirements remain visible
in $\mathcal R(A\otimes\mathbb Q_3)$.
Thus the stated conclusion follows
from the literal coefficient map,
without reversing it.

For a concrete local comparison, the
plus Fermat quartic has no
$\mathbb Q_2$-point. Clear denominators
in a putative point with at least one
odd coordinate: a fourth power is
zero or one modulo $16$, so the sum
of the three coordinates' fourth
powers cannot be zero. Its odd-primary
class-two extensions nevertheless have
the sections constructed in E.17.
Pushing such a section to unipotent
cohomology does not place its
$2$-adic restriction in the image
of an empty local point set. This
is an explicit instance of the
Selmer condition in the same
commuting diagram.

The distinction between full and
finite-support classes is essential
in applying this comparison to the
Fermat tower. E.16 constructs full
pro-prime sections; it makes no
claim that their images satisfy
every Chabauty--Kim point-image
condition. E.9--E.12 instead use
the supplied local theorem for
full profinite sections to obtain
the local geometric conditions and
then prove global cuspidality.
All three conclusions follow along
their displayed maps. The anchors
are Lead 4 and \S51.16, folios
733v--736v, scans 172--178.
\end{proof}

